%------------------------------------------------------------
%these tricks allow to compile the presentation on pdflatex, xelatex and latex+dvipdfm combo.
%However, in the last case the dvi file will be most certainly unreadable
%if you use only one routine, then you can replace this section by:
%
%\documentclass{beamer}
%for xelatex and pdflatex
%
%\documentclass[dvipdfm]{beamer}
%for latex+dvipdfm
%
\RequirePackage{ifpdf}
\RequirePackage{ifxetex}
\ifpdf
  \documentclass{beamer}
\else
  \ifxetex
    \documentclass{beamer}
  \else
    \documentclass[dvipdfm]{beamer}
  \fi
\fi
%-------------------------------------------------------------

%necessary packages
%
\usepackage{pgf,pgfpages}
\usepackage{graphicx}
\usepackage{units}
\usepackage[utf8]{inputenc}

\ifxetex
   \usepackage{fontspec}
\else
   \usepackage[T1]{fontenc} 
\fi  %this trick allows to have good fonts under both pdflatex and xelatex

\newcommand{\donothing}{}

%% end of necessary packages

\usepackage{amsmath, amsthm} %some useful macros for math fonts and environments. Optional


\mode<presentation>
{
  \usetheme{CMLS}
  \setbeamercovered{transparent}
}

\usefonttheme[onlymath]{serif} %optional
\setbeamerfont{frametitle}{size=\Large,series=\bfseries} %setting the font for frametitles


\beamertemplatenavigationsymbolsempty %no navigation symbols - anyway, nobody clicks them. If need this feature, comment/delete this line







\title{Title of the talk}
\author{TZ}
\institute{CMLS, l'X
}
\date{\today}


\begin{document}
%-----------------------------------------------------------------------
% setting the layout for titlepage
%
\setbeamertemplate{footline}[default]
\setbeamertemplate{background}{%
  \vbox to \paperheight{\vfil\hbox to
\paperwidth{\hfil\ifpdf\includegraphics[width=\paperwidth]{images/frontpage-bgp.pdf}
\else\ifxetex\includegraphics[width=\paperwidth]{images/frontpage-bgp.pdf}\else
\includegraphics[width=\paperwidth,natwidth=1261bp,natheight=668bp]{images/frontpage-bgp.pdf}
  \fi
\fi
\hfil}}
}
%-----------------------------------------------------------------------
\begin{frame}
  \titlepage
  \vspace{50pt}
\end{frame}

%------------------------------------------------------------------------
% Set the background for the rest of the slides.
% Insert infoline at the end
%
\setbeamertemplate{background}{\vbox to \paperheight{\hbox to\paperwidth{\hfill\ifpdf\includegraphics[width=\paperwidth]{images/slide-bgp.pdf}
\else\ifxetex\includegraphics[width=\paperwidth]{images/slide-bgp.pdf}\else
\includegraphics[width=\paperwidth,natwidth=1261bp,natheight=668bp]%
{images/slide-bgp.pdf}\fi\fi
\hfil}\vfill}}
\setbeamertemplate{footline}[CMLStheme]
%-------------------------------------------------------------------------
{\setbeamertemplate{footline}{}
\begin{frame}
\frametitle{Table of Contents}
\tableofcontents
\end{frame}
}
\section{Introduction}
\subsection{Basics}
\begin{frame}
	\frametitle{Basics}
Basic functionality of beamer
	
\end{frame}


\subsection{Basic functions}

\begin{frame}
	\frametitle{Formulas}
	Transport equation

	\begin{center}
		 \[\partial_t f+\omega\cdot\nabla_x f =0\]
	\end{center}

\end{frame}

\begin{frame}
\frametitle{Using enumerate}
\begin{enumerate}
\item Lorem 
\item Ipsum
\begin{enumerate}
\item Ipsum-lorem
\item Ipsum-ipsum
\end{enumerate}
\end{enumerate}

\end{frame}

\begin{frame}
\frametitle{Using itemize}
\begin{itemize}
\item Lorem 
\item Ipsum
\begin{itemize}
\item Ipsum-lorem
\item Ipsum-ipsum
\end{itemize}
\end{itemize}
\end{frame}

\begin{frame}
\frametitle{Theorems}
\begin{theorem}
A very important theorem.
\end{theorem}
\begin{proof}
An easy proof.
\end{proof}
\end{frame}


\begin{frame}
\frametitle{Blocks of text}
\begin{block}{Block title}
This is just a block of text
\end{block}

\begin{alertblock}{Alert-block title}
This is a block in alerted colors
\end{alertblock}

\begin{exampleblock}{Example-block title}
This block serves to give examples
\end{exampleblock}
\end{frame}

\subsection{Some tweaks with design}
{\setbeamertemplate{background}{} %no background
\setbeamertemplate{frametitle}{\vspace{0.25cm}\color{Xblue}\insertframetitle} %new colors for the frametitle
\begin{frame}
\frametitle{Frame without background image}


Only infolines are visible.
\end{frame}
}

{\setbeamertemplate{footline}{}% no footer
\begin{frame}
\frametitle{And here we don't have infolines}
\end{frame}
}

{\setbeamertemplate{background}{}%no background
\setbeamertemplate{footline}{}%no footer
\setbeamertemplate{frametitle}{\vspace{0.25cm}\color{Xblue}\insertframetitle}%new colors for the frametitle
\begin{frame}
\frametitle{Pure white frame}

There's nothing here.

\end{frame}
}

\section{Conclusion}
\subsection{One more slide}
\begin{frame}
\frametitle{Last page}
Thank you for reading the whole document!

Feel free to send me any suggestions, remarks, commentaries, requests, etc. at \textit{{\color{Xblue}timofey.zakrevskiy@polytechnique.edu}}
\vfill
Cheers, TZ.
\end{frame}
\end{document}

